\chapter{Título do Apêndice 1}\label{chp:app}

\captionof{listing}{Categorização das Gerações}
\begin{Verbatim}[breaklines=true, frame=single]
ALGORITMO agrupar_geracao
ENTRADA: geracao (lista de valores), limiar (float, opcional = 15)
SAÍDA: grupos (lista de grupos formados), categorias_iniciais (lista de IDs)

INÍCIO
    categorias <- lista vazia
    grupos <- lista vazia
    proximo_id <- 1
    
    // Primeira passagem: agrupamento inicial
    PARA i DE 0 ATÉ tamanho(geracao) - 1:
        valor <- geracao[i]
        
        SE valor <= 20 ENTÃO
            valor <- 0
        
        SE i == 0 ENTÃO
            lim_inf <- valor * (1 - limiar/100)
            lim_sup <- valor * (1 + limiar/100)
            
            grupo <- {
                id: proximo_id,
                media: valor,
                limite_inferior: lim_inf,
                limite_superior: lim_sup,
                valores: [valor]
            }
            
            adicionar grupo em grupos
            adicionar proximo_id em categorias
            proximo_id <- proximo_id + 1
        
        SENÃO
            grupos_validos <- lista vazia
            
            PARA CADA grupo EM grupos:
                SE grupo.limite_inferior <= valor <= grupo.limite_superior ENTÃO
                    SE grupo.media != 0 ENTÃO
                        distancia <- |valor - grupo.media| / grupo.media * 100
                    SENÃO
                        distancia <- |valor| * 100
                    
                    adicionar (grupo, distancia) em grupos_validos
            
            SE grupos_validos NÃO está vazio ENTÃO
                ordenar grupos_validos por distancia (crescente)
                grupo_escolhido <- grupos_validos[0].grupo
                
                adicionar valor em grupo_escolhido.valores
                nova_media <- soma(grupo_escolhido.valores) / tamanho(grupo_escolhido.valores)
                grupo_escolhido.media <- nova_media
                grupo_escolhido.limite_inferior <- nova_media * (1 - limiar/100)
                grupo_escolhido.limite_superior <- nova_media * (1 + limiar/100)
                
                adicionar grupo_escolhido.id em categorias
            
            SENÃO
                lim_inf <- valor * (1 - limiar/100)
                lim_sup <- valor * (1 + limiar/100)
                
                grupo <- {
                    id: proximo_id,
                    media: valor,
                    limite_inferior: lim_inf,
                    limite_superior: lim_sup,
                    valores: [valor]
                }
                
                adicionar grupo em grupos
                adicionar proximo_id em categorias
                proximo_id <- proximo_id + 1
            FIM SE
        FIM SE
    FIM PARA
    
    // Segunda fase: mesclar grupos sobrepostos
    REPETIR
        houve_mesclagem <- FALSO
        grupos_ordenados <- ordenar grupos por media (crescente)
        novos_grupos <- lista vazia
        
        i <- 0
        ENQUANTO i < tamanho(grupos_ordenados):
            grupo_atual <- copia(grupos_ordenados[i])
            j <- i + 1
            
            ENQUANTO j < tamanho(grupos_ordenados):
                proximo_grupo <- grupos_ordenados[j]
                
                SE grupo_atual.limite_inferior <= proximo_grupo.limite_superior E 
                   proximo_grupo.limite_inferior <= grupo_atual.limite_superior ENTÃO
                    
                    adicionar todos valores de proximo_grupo em grupo_atual.valores
                    nova_media <- soma(grupo_atual.valores) / tamanho(grupo_atual.valores)
                    grupo_atual.media <- nova_media
                    grupo_atual.limite_inferior <- nova_media * (1 - limiar/100)
                    grupo_atual.limite_superior <- nova_media * (1 + limiar/100)
                    grupo_atual.id <- minimo(grupo_atual.id, proximo_grupo.id)
                    
                    houve_mesclagem <- VERDADEIRO
                    j <- j + 1
                SENÃO
                    INTERROMPER loop
                FIM SE
            
            adicionar grupo_atual em novos_grupos
            i <- j
        
        grupos <- novos_grupos
    ATÉ houve_mesclagem = FALSO
    
    // Renomear IDs sequencialmente
    grupos_ordenados <- ordenar grupos por media (crescente)
    PARA idx DE 1 ATÉ tamanho(grupos_ordenados):
        grupos_ordenados[idx-1].id <- idx
    
    RETORNAR grupos_ordenados, categorias

FIM ALGORITMO
\end{Verbatim}

\captionof{listing}{Identificação de Transições}
\begin{Verbatim}[breaklines=true, frame=single]
    ALGORITMO identificar_transicoes_e_exportar
ENTRADA: usina (string), geracao (lista de valores), grupos (lista de grupos), 
         categorias_iniciais (lista de IDs), limiar (float, opcional = 15)
SAÍDA: categorias_finais (lista de inteiros)

INÍCIO
    n <- tamanho(geracao)
    
    // Reclassificar todos os valores baseado nos grupos finais
    novas_categorias <- lista vazia
    
    PARA CADA valor_original EM geracao:
        valor <- SE valor_original <= 20 ENTÃO 0 SENÃO valor_original
        
        grupo_encontrado <- NULO
        PARA CADA grupo EM grupos:
            SE grupo.limite_inferior <= valor <= grupo.limite_superior ENTÃO
                grupo_encontrado <- grupo
                INTERROMPER loop
        
        SE grupo_encontrado != NULO ENTÃO
            novo_id <- grupo_encontrado.id
        SENÃO
            grupo_mais_proximo <- grupo em grupos com mínimo |grupo.media - valor|
            novo_id <- grupo_mais_proximo.id
        
        adicionar novo_id em novas_categorias
    
    // Identificar pontos de transição usando janela deslizante
    categorias_finais <- copia(novas_categorias)
    
    PARA i DE 0 ATÉ n-1:
        SE categorias_finais[i] == 0 ENTÃO
            CONTINUAR
        
        // Coletar até 3 categorias anteriores
        categorias_anteriores <- lista vazia
        PARA j DE max(0, i-3) ATÉ i-1:
            adicionar novas_categorias[j] em categorias_anteriores
        
        // Coletar até 3 categorias posteriores
        categorias_posteriores <- lista vazia
        PARA j DE i+1 ATÉ min(n-1, i+3):
            adicionar novas_categorias[j] em categorias_posteriores
        
        SE tamanho(categorias_anteriores) == 0 OU tamanho(categorias_posteriores) == 0 ENTÃO
            CONTINUAR
        
        cat_atual <- novas_categorias[i]
        
        // Contar categorias diferentes
        diferentes_anteriores <- 0
        PARA CADA cat EM categorias_anteriores:
            SE cat != cat_atual ENTÃO
                diferentes_anteriores <- diferentes_anteriores + 1
        
        diferentes_posteriores <- 0
        PARA CADA cat EM categorias_posteriores:
            SE cat != cat_atual ENTÃO
                diferentes_posteriores <- diferentes_posteriores + 1
        
        // Verificar se é ponto de transição
        SE diferentes_anteriores > tamanho(categorias_anteriores)/2 E 
           diferentes_posteriores > tamanho(categorias_posteriores)/2 ENTÃO
            
            valor_atual <- SE geracao[i] <= 20 ENTÃO 0 SENÃO geracao[i]
            
            // Calcular média dos valores anteriores
            valores_anteriores <- lista vazia
            PARA j DE max(0, i-3) ATÉ i-1:
                v <- geracao[j]
                SE v > 20 OU v == 0 ENTÃO
                    adicionar v em valores_anteriores
            
            SE valores_anteriores NÃO vazio ENTÃO
                media_anterior <- soma(valores_anteriores) / tamanho(valores_anteriores)
            SENÃO
                media_anterior <- valor_atual
            
            // Calcular média dos valores posteriores
            valores_posteriores <- lista vazia
            PARA j DE i+1 ATÉ min(n-1, i+3):
                v <- geracao[j]
                SE v > 20 OU v == 0 ENTÃO
                    adicionar v em valores_posteriores
            
            SE valores_posteriores NÃO vazio ENTÃO
                media_posterior <- soma(valores_posteriores) / tamanho(valores_posteriores)
            SENÃO
                media_posterior <- valor_atual
            
            // Calcular diferenças percentuais
            diff_anterior <- |valor_atual - media_anterior| / max(|media_anterior|, 1) * 100
            diff_posterior <- |valor_atual - media_posterior| / max(|media_posterior|, 1) * 100
            
            SE diff_anterior > limiar E diff_posterior > limiar ENTÃO
                categorias_finais[i] <- 0
    
    // Criar diretório de saída
    dir <- caminho_join('Usinas', usina, 'Infos de Geração')
    SE diretorio NÃO existe ENTÃO
        criar_diretorio(dir)
    
    // Gerar arquivo com informações dos grupos
    dados_grupos <- lista vazia
    PARA CADA grupo EM grupos (ordenados por media crescente):
        qtd <- 0
        PARA i DE 0 ATÉ n-1:
            SE categorias_finais[i] == grupo.id ENTÃO
                qtd <- qtd + 1
        
        adicionar em dados_grupos: {
            'Grupo': grupo.id,
            'Média': arredondar(grupo.media, 2),
            'Limite Inferior': arredondar(grupo.limite_inferior, 2),
            'Limite Superior': arredondar(grupo.limite_superior, 2),
            'Registros': qtd
        }
    
    salvar_parquet(dados_grupos, caminho_join(dir, 'Grupos de Geração.parquet'))
    
    // Gerar arquivo com os pontos de transição
    dados_transicao <- lista vazia
    PARA i DE 0 ATÉ n-1:
        SE categorias_finais[i] == 0 ENTÃO
            adicionar em dados_transicao: {
                'Índice': i,
                'Geração': arredondar(geracao[i], 3)
            }
    
    salvar_parquet(dados_transicao, caminho_join(dir, 'Transições de Geração.parquet'))
    
    RETORNAR categorias_finais

FIM ALGORITMO
\end{Verbatim}

\captionof{listing}{Modelo AMPL}
\begin{Verbatim}[breaklines=true, frame=single]
ALGORITMO modelo_regressao_emissoes

// =================
// DECLARAÇÃO DE CONJUNTOS
// =================

CONJUNTO ANO
CONJUNTO HORARIO[a] para cada a em ANO, ordenado: 0 até H[a]

// =================
// DECLARAÇÃO DE PARÂMETROS
// =================

PARÂMETRO H[ANO]                  // Número de horas por ano
PARÂMETRO emissoes[ANO]           // Emissões anuais em kg CO2e
PARÂMETRO pg[ANO, HORARIO]        // Geração por ano e hora

PARÂMETRO S_base <- 100           // Base para escala (não utilizado explicitamente)
PARÂMETRO lambda <- 30            // Parâmetro de regularização L2

// =================
// DECLARAÇÃO DE VARIÁVEIS
// =================

// Variáveis para regressão padrão
VAR Alpha >= 0                     // Coeficiente quadrático (cargas muito altas)
VAR Beta                           // Coeficiente linear (eficiência em carga parcial)
VAR Gamma                          // Termo constante
VAR Omega                          // Amplitude do termo exponencial
VAR Mu                             // Taxa de crescimento exponencial

// Variáveis para regressão com regularização L2
VAR Alpha_L2 >= 0
VAR Beta_L2
VAR Gamma_L2
VAR Omega_L2
VAR Mu_L2

// Variáveis derivadas (emissões horárias estimadas)
VAR Emissao_horaria[ANO, HORARIO] =
    Alpha * pg[a,h]^2
    + Beta * pg[a,h]
    + Gamma
    + Omega * exp(Mu * pg[a,h])

VAR Emissao_horaria_L2[ANO, HORARIO] =
    Alpha_L2 * pg[a,h]^2
    + Beta_L2 * pg[a,h]
    + Gamma_L2
    + Omega_L2 * exp(Mu_L2 * pg[a,h])

// =================
// RESTRIÇÕES FÍSICAS
// =================

// Restrição 1: Emissões horárias nunca negativas
PARA CADA a em ANO:
    PARA CADA h em HORARIO[a]:
        SUJEITO A Emissao_nao_negativa[a,h]:
            Alpha * pg[a,h]^2
            + Beta * pg[a,h]
            + Gamma
            + Omega * exp(Mu * pg[a,h]) >= 0

        SUJEITO A Emissao_nao_negativa_L2[a,h]:
            Alpha_L2 * pg[a,h]^2
            + Beta_L2 * pg[a,h]
            + Gamma_L2
            + Omega_L2 * exp(Mu_L2 * pg[a,h]) >= 0

// Restrição 2: Emissão zero quando geração é zero
SUJEITO A Emissao_no_zero:
    Gamma + Omega >= 0

SUJEITO A Emissao_no_zero_L2:
    Gamma_L2 + Omega_L2 >= 0

// =================
// FUNÇÕES OBJETIVO
// =================

// Objetivo 1: MSE (Erro Quadrático Médio) - Regressão padrão
FUNCAO_OBJETIVO MSE:
    MINIMIZAR (1 / SOMA_{a em ANO} |HORARIO[a]|) 
    * SOMA_{a em ANO} (SOMA_{h em HORARIO[a]} Emissao_horaria[a,h] - emissoes[a])^2

// Objetivo 2: MSE com regularização L2
FUNCAO_OBJETIVO MSE_L2:
    MINIMIZAR (1 / SOMA_{a em ANO} |HORARIO[a]|)
    * SOMA_{a em ANO} (SOMA_{h em HORARIO[a]} Emissao_horaria_L2[a,h] - emissoes[a])^2
    + lambda * (Alpha_L2^2 + Beta_L2^2 + Gamma_L2^2 + Omega_L2^2 + Mu_L2^2)

// =================
// DEFINIÇÃO DOS PROBLEMAS
// =================

PROBLEMA Regressao_Estatica:
    FUNCAO_OBJETIVO: MSE
    VARIÁVEIS: Emissao_horaria, Alpha, Beta, Gamma, Omega, Mu
    RESTRIÇÕES: Emissao_nao_negativa, Emissao_no_zero

PROBLEMA Regressao_Estatica_L2:
    FUNCAO_OBJETIVO: MSE_L2
    VARIÁVEIS: Emissao_horaria_L2, Alpha_L2, Beta_L2, Gamma_L2, Omega_L2, Mu_L2
    RESTRIÇÕES: Emissao_nao_negativa_L2, Emissao_no_zero_L2

FIM ALGORITMO
\end{Verbatim}

\captionof{listing}{Geração das Emissões Sintéticas}
\begin{Verbatim}[breaklines=true, frame=single]
ALGORITMO geracao_sinteticas
SAÍDA: Arquivos Parquet salvos em disco (sem retorno explícito)

INÍCIO

    // ==================== INICIALIZAÇÃO ====================
    
    horas <- dicionário vazio
    anos <- listar arquivos no diretório 'IEMA/Tabelas'
    remover 'desktop.ini' do conjunto
    para cada arquivo, remover extensão (ex: '.parquet')
    ordenar anos em ordem crescente
    
    // Calcular horas por ano (considerando anos bissextos)
    PARA CADA ano EM anos:
        converter ano para inteiro
        SE ano % 4 == 0 ENTÃO
            horas[ano] <- 8784    // Ano bissexto (366 dias)
        SENÃO
            horas[ano] <- 8760    // Ano normal (365 dias)
        FIM SE
    FIM PARA
    
    // ==================== PROCESSAMENTO POR USINA ====================
    
    PARA CADA usina EM listar_diretorios('Usinas'):
        
        // Carregar dados de geração da usina
        caminho_pg <- caminho('Usinas', usina, 'Dados Externos', 'ONS.parquet')
        PG <- ler_parquet(caminho_pg)['Geração']  // Array/Series de geração
        
        pastas <- ['Padrão', 'Regressão L2']
        
        // ==================== CÁLCULO DE EMISSÕES SINTÉTICAS ====================
        
        PARA CADA pasta EM pastas:
            
            caminho_coeficientes <- caminho('Usinas', usina, 'Minimização', 
                                           pasta, 'Emissões', 'Coeficientes.parquet')
            
            SE arquivo_existe(caminho_coeficientes) ENTÃO
                
                // Carregar tabela de coeficientes
                tab <- ler_parquet(caminho_coeficientes)
                
                // Extrair coeficientes da equação
                alpha <- valor de tab onde Coeficientes == 'Alpha'
                beta  <- valor de tab onde Coeficientes == 'Beta'
                gamma <- valor de tab onde Coeficientes == 'Gamma'
                omega <- valor de tab onde Coeficientes == 'Omega'
                mu    <- valor de tab onde Coeficientes == 'Mu'
                
                // Calcular emissões horárias sintéticas
                // Equação: emissão = $alpha$*(PG/100)² + $beta$*(PG/100) + $gamma$ + $omega$*e^($mu$*(PG/100))
                emissao <- lista vazia
                
                PARA k DE 0 ATÉ comprimento(PG) - 1:
                    pg_normalizado <- PG[k] / 100
                    
                    termo_quadratico <- alpha * (pg_normalizado^2)
                    termo_linear    <- beta * pg_normalizado
                    termo_constante <- gamma
                    termo_exponencial <- omega * exp(mu * pg_normalizado)
                    
                    valor_emissao <- termo_quadratico + termo_linear + 
                                    termo_constante + termo_exponencial
                    
                    adicionar arredondar(valor_emissao, 3) em emissao
                FIM PARA
                
                // ==================== CRIAÇÃO DO DATAFRAME HORÁRIO ====================
                
                horas_expandidas <- lista vazia
                anos_expandidos <- lista vazia
                
                // Expandir horas e anos para cada período
                PARA CADA ano EM anos:
                    inicio <- 0
                    num_horas <- horas[ano]
                    
                    // Adicionar índices de 0 até num_horas-1
                    horas_expandidas.extend(lista de 0 até num_horas-1)
                    
                    // Adicionar o ano repetido num_horas vezes
                    anos_expandidos.extend([ano] * num_horas)
                FIM PARA
                
                // Criar DataFrame com dados horários
                df_horario <- DataFrame({
                    'Ano': anos_expandidos,
                    'Índice': horas_expandidas,
                    'Emissão': emissao
                })
                
                // Salvar arquivo de emissões horárias
                caminho_horario <- caminho('Usinas', usina, 'Minimização', 
                                          pasta, 'Emissões', 'Horárias.parquet')
                salvar_parquet(df_horario, caminho_horario, index=FALSO)
                
                // ==================== CÁLCULO DE EMISSÕES ANUAIS ====================
                
                // Agrupar por ano e somar emissões
                df_anual <- df_horario.agrupar_por('Ano')['Emissão'].somar()
                df_anual <- resetar_index(df_anual)
                
                // Carregar emissões reais do IEMA
                caminho_iema <- caminho('Usinas', usina, 'Dados Externos', 'IEMA.parquet')
                iema <- ler_parquet(caminho_iema)['Emissões']
                emissoes_reais <- converter_para_lista(iema)
                
                // Obter emissões anuais sintéticas
                emissoes_sinteticas <- df_anual['Emissão'].converter_para_lista()
                
                // Calcular K_i = Emissão_Sintética / Emissão_Real
                valores_ki <- lista vazia
                
                PARA j DE 0 ATÉ comprimento(emissoes_sinteticas) - 1:
                    ki <- emissoes_sinteticas[j] / emissoes_reais[j]
                    adicionar arredondar(ki, 3) em valores_ki
                FIM PARA
                
                // Adicionar coluna K_i ao DataFrame anual
                df_anual['K_i'] <- valores_ki
                
                // Salvar arquivo de emissões anuais
                caminho_anual <- caminho('Usinas', usina, 'Minimização', 
                                        pasta, 'Emissões', 'Anuais.parquet')
                salvar_parquet(df_anual, caminho_anual, index=FALSO)
                
            FIM SE
            
        FIM PARA
        
    FIM PARA

FIM ALGORITMO
\end{Verbatim}

\captionof{listing}{Preparação das \emph{Features} da Rede Neural}
\begin{Verbatim}[breaklines=true, frame=single]
ALGORITMO tratamento_dados
ENTRADA: iema_recente (DataFrame)
SAÍDA: nenhum retorno explícito (salva arquivos em disco)

INÍCIO
    // Definir pastas a processar baseado no tipo de previsão
    pastas <- ['Padrão', 'Regressão L2']
    
    // Processar cada usina
    PARA CADA usina EM listar_diretorios('Usinas'):
        
        PARA CADA pasta EM pastas:
            
            // Carregar dados de geração e emissão
            geracao <- ler_parquet(caminho('Usinas', usina, 'Dados Externos', 'ONS.parquet'))
            emissao <- ler_parquet(caminho('Usinas', usina, 'Minimização', pasta, 'Emissões', 'Horárias.parquet'))
            
            // Carregar dados unitários (tentando por 'Usina' ou 'CEG')
            TENTAR
                unitarios <- iema_recente[iema_recente['Usina'] == usina]
                remover coluna 'Usina' de unitarios
                resetar índice de unitarios
            EXCETO KeyError
                unitarios <- iema_recente[iema_recente['CEG'] == usina]
                remover coluna 'CEG' de unitarios
                resetar índice de unitarios
            FIM TENTAR
            
            // Unificar os dados
            df <- mesclar(geracao, emissao, how='left')
            df <- mesclar(df, unitarios, how='cross')
            
            // == TRANSFORMAÇÃO DAS HORAS EM VALORES CÍCLICOS ==
            horas <- dicionário vazio
            anos <- listar arquivos em 'IEMA/Tabelas' (excluindo 'desktop.ini')
            remover extensão '.parquet' dos nomes
            ordenar anos
            
            PARA CADA ano EM anos:
                SE ano % 4 == 0 ENTÃO
                    horas[ano] <- 8784
                SENÃO
                    horas[ano] <- 8760
                FIM SE
            
            H <- mapear df['Ano'] usando dicionário horas
            theta <- 2 * PI * df['Índice'] / H
            
            df['Cosseno Índice'] <- cos(theta)
            df['Seno Índice'] <- sin(theta)
            
            // == FEATURES DE TRANSIÇÃO ==
            // Duração da transição
            transicao_list <- converter (df['Categoria de geração'] == 0) para inteiro
            duracao <- vetor de zeros (mesmo tamanho de transicao_list)
            contador <- 0
            
            PARA i DE 0 ATÉ tamanho(transicao_list) - 1:
                SE transicao_list[i] == 1 ENTÃO
                    contador <- contador + 1
                SENÃO
                    contador <- 0
                duracao[i] <- contador
            df['Duração da Transição'] <- duracao
            
            // Fase da transição
            fase <- vetor de zeros (mesmo tamanho de transicao_list)
            
            PARA i DE 0 ATÉ tamanho(transicao_list) - 1:
                SE transicao_list[i] == 1 ENTÃO  // em transição
                    inicio <- (i == 0) OU (transicao_list[i-1] == 0)
                    fim <- (i == tamanho(transicao_list)-1) OU (transicao_list[i+1] == 0)
                    
                    SE inicio E fim ENTÃO
                        fase[i] <- 1  // transição de uma hora
                    SENÃO SE inicio ENTÃO
                        fase[i] <- 1  // início
                    SENÃO SE fim ENTÃO
                        fase[i] <- 3  // fim
                    SENÃO
                        fase[i] <- 2  // meio
                    FIM SE
                FIM SE
            df['Fase da Transição'] <- fase
            
            // Separar dados em treino, validação e teste
            removidas <- ['Delta menos', 'Delta mais', 'Índice']
            adicionar 'Ano' a removidas
            FIM SE
            
            ano_divisao <- 2023
            
            x_tr <- df[df['Ano'] < ano_divisao]
            resetar índice de x_tr
            remover colunas removidas de x_tr
            
            x_val <- df[df['Ano'] == ano_divisao]
            resetar índice de x_val
            remover colunas removidas de x_val
            
            x_te <- df[df['Ano'] > ano_divisao]
            resetar índice de x_te
            remover colunas removidas de x_te
            
            // Definir colunas categóricas e numéricas
            categoricas <- ['Categoria de geração']
            numericas <- ['Geração', 'Duração da Transição', 'Fase da Transição']
            
            // Configurar transformador de colunas
            transf <- ColumnTransformer([
                ('onehot', OneHotEncoder(sparse_output=FALSO, handle_unknown='ignore'), categoricas),
                ('scaler', StandardScaler(), numericas)
            ], remainder='passthrough', verbose_feature_names_out=FALSO)
            
            // Aplicar transformações
            data_tr <- transf.fit_transform(x_tr)
            data_val <- transf.transform(x_val)
            data_te <- transf.transform(x_te)
            
            // Criar DataFrames finais
            df_tr <- DataFrame(data_tr, colunas = transf.get_feature_names_out())
            df_val <- DataFrame(data_val, colunas = transf.get_feature_names_out())
            df_te <- DataFrame(data_te, colunas = transf.get_feature_names_out())
            
            // Extrair scaler para coluna 'Geração'
            scaler <- transf.named_transformers_['scaler']
            idx_pg <- posição de 'Geração' em numericas
            
            df_norm <- DataFrame({
                'Variável': ['Geração'],
                'Mean': [scaler.mean_[idx_pg]],
                'Std': [scaler.scale_[idx_pg]]
            })
            
            // Definir diretório de saída
            dir <- caminho('Usinas', usina, 'Minimização', pasta, 'Rede Neural', 'Dados')
            
            // Criar diretório se não existir
            criar_diretorio(dir, exist_ok=VERDADEIRO)
            
            // Salvar DataFrames
            salvar_parquet(df_tr, caminho(dir, 'Treino.parquet'), index=FALSO)
            salvar_parquet(df_val, caminho(dir, 'Validação.parquet'), index=FALSO)
            salvar_parquet(df_te, caminho(dir, 'Teste.parquet'), index=FALSO)
            salvar_parquet(df_norm, caminho(dir, 'Scaler Geração.parquet'), index=FALSO)
            
        FIM PARA
    FIM PARA

FIM ALGORITMO
\end{Verbatim}

\captionof{listing}{Treinamento da Rede Neural}
\begin{Verbatim}[breaklines=true, frame=single]
ALGORITMO treinar_modelo
ENTRADA: 
    model (rede neural)
    train_loader (dados de treino)
    val_loader (dados de validação)
    dyn_cols (colunas dinâmicas)
    epochs (inteiro, padrão = 100)
    patience (inteiro, padrão = 50)
    huber_delta (float, padrão = 0.5)
    peso_transicao (float, padrão = 8.0)
    clip_grad (float, padrão = 1.0)

SAÍDA:
    model (modelo treinado)
    history (dicionário com histórico de métricas)

INÍCIO
    // Mover modelo para dispositivo (CPU/GPU)
    model <- mover_para_device(model, device)
    
    // Configurar otimizador e funções de perda
    optimizer <- Adam(model.parameters(), lr=0.001)
    criterion_huber <- HuberLoss(delta=huber_delta)
    criterion_mae <- L1Loss()
    criterion_mse <- MSELoss()
    
    // Configurar scheduler para reduzir learning rate
    scheduler <- ReduceLROnPlateau(optimizer, mode='min', factor=0.5, 
                                   patience=5, min_lr=0.00001)
    
    // Inicializar variáveis de controle
    best_val_loss <- infinito
    patience_counter <- 0
    best_model_state <- None
    history <- {
        'loss': [], 'mae': [], 'mse': [],
        'val_loss': [], 'val_mae': [], 'val_mse': []
    }
    
    // Loop de épocas
    PARA epoch DE 1 ATÉ epochs:
        
        // === FASE DE TREINAMENTO ===
        model.treino()
        train_loss <- 0.0
        train_mae <- 0.0
        train_mse <- 0.0
        
        PARA CADA (x_dyn, x_sta, y) EM train_loader:
            // Mover dados para o dispositivo
            x_dyn <- mover_para_device(x_dyn, device)
            x_sta <- mover_para_device(x_sta, device)
            y <- mover_para_device(y, device)
            y <- squeeze(y)  // Remover dimensões extras
            
            // Zerar gradientes
            optimizer.zerar_gradientes()
            
            // Forward pass
            output <- model(x_dyn, x_sta)
            
            // Calcular perda base (Huber)
            loss_data <- criterion_huber(output, y)
            
            // Aplicar pesos para transições
            is_transicao <- (x_sta[:, 0] == 1)
            weight <- onde(is_transicao, peso_transicao, 1.0)
            loss <- media(loss_data * weight)
            
            // Backward pass
            loss.backward()
            
            // Clip gradients para evitar explosão
            clip_grad_norm_(model.parameters(), max_norm=clip_grad)
            
            // Atualizar pesos
            optimizer.step()
            
            // Acumular métricas
            train_loss <- train_loss + loss.item()
            train_mae <- train_mae + media(criterion_mae(output, y)).item()
            train_mse <- train_mse + media(criterion_mse(output, y)).item()
        
        // Calcular médias das métricas de treino
        train_loss <- train_loss / tamanho(train_loader)
        train_mae <- train_mae / tamanho(train_loader)
        train_mse <- train_mse / tamanho(train_loader)
        
        // === FASE DE VALIDAÇÃO ===
        model.avaliacao()  // Modo de avaliação (sem dropout, etc.)
        val_loss <- 0.0
        val_mae <- 0.0
        val_mse <- 0.0
        
        // Desabilitar cálculo de gradientes para validação
        COM torch.no_grad():
            PARA CADA (x_dyn, x_sta, y) EM val_loader:
                x_dyn <- mover_para_device(x_dyn, device)
                x_sta <- mover_para_device(x_sta, device)
                y <- mover_para_device(y, device)
                y <- squeeze(y)
                
                output <- model(x_dyn, x_sta)
                
                loss_val <- media(criterion_huber(output, y))
                val_loss <- val_loss + loss_val.item()
                val_mae <- val_mae + criterion_mae(output, y).item()
                val_mse <- val_mse + criterion_mse(output, y).item()
        
        // Calcular médias das métricas de validação
        val_loss <- val_loss / tamanho(val_loader)
        val_mae <- val_mae / tamanho(val_loader)
        val_mse <- val_mse / tamanho(val_loader)
        
        // Registrar histórico
        adicionar train_loss em history['loss']
        adicionar train_mae em history['mae']
        adicionar train_mse em history['mse']
        adicionar val_loss em history['val_loss']
        adicionar val_mae em history['val_mae']
        adicionar val_mse em history['val_mse']
        
        // Atualizar scheduler baseado na loss de validação
        scheduler.step(val_loss)
        
        // Verificar se é o melhor modelo até agora
        SE val_loss < best_val_loss ENTÃO
            best_val_loss <- val_loss
            patience_counter <- 0
            best_model_state <- copiar(model.state_dict())
        SENÃO
            patience_counter <- patience_counter + 1
            
            SE patience_counter >= patience ENTÃO
                INTERROMPER loop  // Early stopping
            FIM SE
        FIM SE
        
    FIM PARA
    
    // Carregar o melhor modelo encontrado
    SE best_model_state NÃO é None ENTÃO
        model.load_state_dict(best_model_state)
    FIM SE
    
    RETORNAR model, history

FIM ALGORITMO
\end{Verbatim}